\clearpage
\VerdierPhysicalPageAnchor{0161}%
\begingroup\small\setlength{\parskip}{0pt}\setlength{\emergencystretch}{1.5em}
\VerdierSemanticAnchor{ch3-3-2-9-exactness-test}
$\mathcal D^\wedge$ で完全である。(3.2.9.1) が完全であることは、
$\mathcal D$ の任意の対象 $W$ に対して
\[
\operatorname{Hom}_{\mathcal D^\wedge}(h(W),h(X))\longrightarrow
\operatorname{Hom}_{\mathcal D^\wedge}(h(W),h(Y))\longrightarrow
\operatorname{Hom}_{\mathcal D^\wedge}(h(W),h(Z))
\]
が完全であることと同値である。これは $h$ の充満忠実性と (1.2.1) から
従う。

\VerdierSemanticAnchor{ch3-3-2-9-factorization}
$\mathcal A$ をアーベル圏とする。$\mathcal H$ との合成をとる関手
\[
\mathcal E x(\mathrm A(\mathcal D),\mathcal A)
\longrightarrow\mathcal H\!om(\mathcal D,\mathcal A)
\]
は明らかに充満忠実である。実際、完全関手間の射は、十分多くの射影対象上での値を
知れば一意に定まる。任意のコホモロジー的関手
$\Psi\colon\mathcal D\to\mathcal A$ を $\mathrm A(\mathcal D)$ を経由して
分解することを示せばよい：
\VerdierSemanticAnchor{ch3-3-2-9-2}
\[
\begin{tikzcd}[column sep=5em,row sep=3em]
\mathcal D\arrow[r,"\mathcal H"]\arrow[dr,"\Psi"']&
\mathrm A(\mathcal D)\arrow[d,"\Phi"]\\
&\mathcal A
\end{tikzcd}
\tag{3.2.9.2}
\]
ここで $\Phi$ は完全である。$\mathrm A(\mathcal D)$ の各対象
$X\xrightarrow fY$ に対し
\begingroup\renewcommand{\thefootnote}{(*)}%
\[
\Phi(X\xrightarrow fY)=
\operatorname{Image}_{\mathcal A}
(\Psi(X)\xrightarrow{\Psi(f)}\Psi(Y)).\footnote{関手 $\Phi$ のこの構成では、
$\mathcal A$ における「像」関手を選ぶ必要がある。この選択はもちろん、
一意な同型を除いて一意である。また、以後そうすることにするが、この
「像」関手が $\mathcal A$ の各恒等射 $\operatorname{id}_T$ に、射
$T\xrightarrow{\operatorname{id}_T}T$ を備えた対象 $T$ を対応させるよう
課してよい。}
\]
\endgroup
と定める。これにより、図式 (3.2.9.2) を可換にする
$\mathrm A(\mathcal D)$ から $\mathcal A$ への関手 $\Phi$ が定まる。$\Phi$ が完全で
あることを示そう。(3.2.7) から $\Phi$ は直ちに左完全であり、反対圏を取れば、
同様に右完全であることが分かる。これで定理 (3.2.1) の証明が完了する。

\VerdierSemanticAnchor{ch3-3-2-10}
\medskip\noindent\textbf{注意 3.2.10.}\quad
1) 定理 (3.2.1, c)) に記された普遍性は、圏同値を除いて
$\mathrm A(\mathcal D)$ を決定する。

2) $\mathcal D$ で可算直和または可算直積が表現可能なら、
$\mathrm A(\mathcal D)$ のすべての射影対象（それぞれ入射対象）は
$\mathcal D$ の対象と同型である。実際、そのような対象は
$\mathrm A(\mathcal D)$ において $\mathcal D$ のある対象の直和因子であり、
従って $\mathcal D$ のある対象上の射影子の核である。
上の仮定のもとでは、射影子は
$\mathcal D$ において分解可能である (1.2.9)。
\endgroup
